% Chapter 2: Your First CLAUDE.md
\chapter{Your First CLAUDE.md}
\label{ch:first-claude-md}

\begin{itshape}
Claude keeps asking you the same questions about your project.
It does not know your build command, your test runner, or your architecture.
Every session starts with you re-explaining context that should be obvious.
The fix takes five minutes.
\end{itshape}

\section{What CLAUDE.md Does}
\label{sec:what-claude-md}

\marginnote[Official]{Run \code{/init} to auto-generate a starter CLAUDE.md
from your project structure.\sourceurl{https://docs.anthropic.com/en/docs/claude-code/best-practices}}

\official{The \code{CLAUDE.md} file is the single most important configuration artifact.}
Claude reads it at session start and uses it to understand your project's conventions,
build commands, architecture decisions, and rules of engagement.

Think of it as a briefing document for a new team member who is fast,
capable, and amnesiac---they forget everything between sessions
unless you write it down.

\subsection{Hierarchical Loading}

Claude loads \code{CLAUDE.md} files from multiple locations, merging them in order:

\begin{enumerate}
  \item \code{\textasciitilde/.claude/CLAUDE.md} --- global defaults (applies to all projects)
  \item Project root \code{CLAUDE.md} or \code{.claude/CLAUDE.md} --- project-level instructions
  \item Parent and child directory \code{CLAUDE.md} files --- subdirectory overrides
\end{enumerate}

For now, you only need the project-level file. The global file and
advanced architecture come later (\cref{ch:architecture}).

% -------------------------------------------------------------------
\section{Generate It: \texttt{/init}}
\label{sec:init}

The fastest path to a useful CLAUDE.md:

\begin{enumerate}
  \item Open your project in Claude Code.
  \item Type \code{/init}.
  \item Claude examines your project structure and proposes a starter CLAUDE.md.
  \item Review, edit, commit.
\end{enumerate}

The auto-generated version is a solid starting point.
What follows are the refinements that make it effective.

% -------------------------------------------------------------------
\section{What to Include}
\label{sec:what-to-include}

\convergence{Include what Claude cannot infer. Exclude what it can.}

\begin{description}
  \item[Build and test commands] \code{make pilot}, \code{pytest tests/},
    \code{npm run build}---anything non-standard or project-specific.
  \item[Code style rules] Only those that differ from standard conventions.
    Do not restate PEP~8 or Prettier defaults.
  \item[Architecture decisions] Where things go and why.
    ``API routes in \code{src/routes/}. Each route has a test in \code{tests/routes/}.''
  \item[Verification criteria] The single highest-leverage inclusion.
    ``Always run tests after code changes. Never commit without passing CI.''
\end{description}

\begin{fullwidth}
\begin{beforeafter}
\before
\begin{minted}{markdown}
# My Project
This is a web app built with React and Express.
We use TypeScript and try to follow best practices.
Tests are important to us.
\end{minted}

\after
\begin{minted}{markdown}
# My Project

## Build
- `npm run build` — compile TypeScript
- `npm test` — run Jest tests
- `npm run lint` — ESLint + Prettier

## Architecture
- src/routes/ — API endpoints (one file per resource)
- src/services/ — Business logic (no HTTP concerns)
- tests/ — mirrors src/ structure

## Standards
- TypeScript strict mode, no `any`
- Run tests after every code change
- Never commit without passing lint + tests
\end{minted}
\end{beforeafter}
\end{fullwidth}

\margintip{The ``before'' tells Claude about your project.
The ``after'' tells Claude how to \emph{work} on your project.
Commands, locations, and rules---not descriptions.}

% -------------------------------------------------------------------
\section{What to Exclude}
\label{sec:what-to-exclude}

\begin{itemize}
  \item What Claude can infer from the codebase (standard language conventions)
  \item Long explanations---keep instructions terse and actionable
  \item Aspirational rules nobody follows---document what IS, not what you wish
  \item Information that changes frequently (put it in memory instead---see \cref{ch:context})
\end{itemize}

\begin{tipbox}[Quality Test]
If Claude keeps violating a rule, one of three things is true:
(1) the CLAUDE.md is too long and the rule is getting lost,
(2) the rule contradicts another rule, or
(3) the rule should be a hook instead of an instruction (see \cref{ch:extending}).
\end{tipbox}

% -------------------------------------------------------------------
\section{Your First Deny Rules}
\label{sec:first-deny}

\marginnote[Official]{Settings documentation:\sourceurl{https://docs.anthropic.com/en/docs/claude-code/settings}}

Before Claude touches any code, protect your secrets.
Create \code{.claude/settings.json}:

\begin{minted}{json}
{
  "permissions": {
    "deny": [
      "Read(.env*)",
      "Read(secrets/**)",
      "Read(**/*.pem)",
      "Read(**/*.key)",
      "Read(**/credentials*)"
    ]
  }
}
\end{minted}

\margintip{Deny rules are absolute. Unlike CLAUDE.md instructions (which are advisory),
deny rules cannot be overridden by conversation or prompt injection.}

Deny rules are the one configuration that is never wrong to add on day one.
They cost nothing in context, they never cause false negatives in normal work,
and they prevent the one class of mistake that is truly irreversible:
leaking secrets into a commit, a log, or a conversation.

% -------------------------------------------------------------------
\section{Quick Reference}

\begin{itemize}
  \item Run \code{/init} to generate a starter CLAUDE.md
  \item Include: build commands, architecture, code style, verification criteria
  \item Exclude: what Claude can infer, long prose, aspirational rules
  \item Add deny rules for secrets in \code{.claude/settings.json} on day one
  \item Commit both files to version control
\end{itemize}

% -------------------------------------------------------------------
\begin{trybox}
\begin{enumerate}
  \item Run \code{/init} on your current project.
  \item Compare the generated CLAUDE.md to what you actually need.
    Add build/test commands if missing.
  \item Create \code{.claude/settings.json} with deny rules for \code{.env*}
    and any secrets directories.
  \item Start a new session and verify Claude reads your CLAUDE.md
    (ask: ``What are my build commands?'').
\end{enumerate}
Time: $\sim$5 minutes.
\end{trybox}
